\chapter{Diagramas de Bloques de Confiabilidad Coloreados}
\label{chapter:crbd}
Una vez estudiado el modelo fundamental de la teoría y sus problemáticas subyacentes, en este capítulo se definirá una extensión del mismo, los RBD coloreados (cRBD), que solucionan el problema de la composición paralela y, no menos importante, garantizan propiedades que se desean cumplir sobre el modelo: la composicionalidad, la modularidad y la clausura bajo operaciones. Los cRBD se establecerán en términos de RBD, con la diferencia que el grafo podrá tener ciclos

\section{cRBD: Definición formal}
\label{chapter:crbd_definicion}
\completar{Definición de cRBD sobre grafos dirigidos (permitiendo ciclos/dependencias mutuas), alfabeto de colores y proyecciones.}

\section{Semántica}
\label{sec:crbd_semantica}
\completar{Semántica multicapa, compartición física de nodos entre distintos colores y condición de éxito del sistema.}

\section{Operadores}
\label{sec:crbd_operadores}
\completar{Operaciones de composición en cRBDs y acoplamiento entre canales.}